%======================================================================
% SANDRA ECOSYSTEM - Diseño Revista Premium
%======================================================================

\documentclass[11pt, spanish, letterpaper, twocolumn, openany]{book}

\usepackage{fontspec}
\setmainfont{Inter}[Path = assets/fonts/, Extension = .otf, UprightFont = *-Regular, BoldFont = *-Bold]

\usepackage[
    top=1.5cm,
    bottom=1.8cm,
    left=1.5cm,
    right=1.5cm,
    columnsep=14pt
]{geometry}

\usepackage{graphicx}
\usepackage{fancyhdr}
\usepackage{titlesec}
\usepackage{xcolor}
\usepackage{microtype}
\usepackage{tikz}
\usetikzlibrary{calc}
\usepackage[most]{tcolorbox}
\usepackage{parskip}
\usepackage{booktabs}
\usepackage{setspace}
\usepackage{changepage}

%--------------------------------------------------------------------
% COLORES
%--------------------------------------------------------------------
\definecolor{sandraDark}{HTML}{1B2631}
\definecolor{sandraTeal}{HTML}{00796B}
\definecolor{sandraGold}{HTML}{FF8F00}
\definecolor{sandraLight}{HTML}{E0F2F1}
\definecolor{sandraCream}{HTML}{F5F5F5}
\definecolor{sandraGray}{HTML}{78909C}
\definecolor{isoBlue}{HTML}{1565C0}
\definecolor{isoOrange}{HTML}{EF6C00}
\definecolor{isoAmber}{HTML}{FFF8E1}

%--------------------------------------------------------------------
% ESPACIADO
%--------------------------------------------------------------------
\setstretch{1.45}
\setlength{\parskip}{0.8ex}

%--------------------------------------------------------------------
% CABECERA
%--------------------------------------------------------------------
\pagestyle{fancy}
\fancyhf{}
\fancyhead[RO]{\color{sandraTeal}\fontsize{9}{11}\selectfont\scshape\rightmark}
\fancyhead[LE]{\color{sandraTeal}\fontsize{9}{11}\selectfont\scshape\leftmark}
\fancyfoot[CO]{\color{sandraGray}\thepage}
\fancyfoot[RE,LO]{\color{sandraGray}\fontsize{7}{9}\selectfont\itshape Sandra Ecosystem}

%--------------------------------------------------------------------
% TÍTULOS
%--------------------------------------------------------------------
\titleformat{\section}{\normalfont\large\bfseries\color{sandraDark}}{\thesection}{0.5em}{}
\titlespacing*{\section}{0pt}{1.8ex}{0.8ex}

%--------------------------------------------------------------------
% CAJAS
%--------------------------------------------------------------------
\newtcolorbox{notaeditorial}{
    colback=sandraCream,
    colframe=sandraTeal,
    boxrule=0pt,
    left=8pt,
    right=8pt,
    top=8pt,
    bottom=8pt,
    borderline west={2pt}{0pt}{sandraTeal},
    before skip=10pt,
    after skip=10pt
}

\newtcolorbox{estandarISO}{
    colback=sandraLight,
    colframe=isoBlue,
    boxrule=1pt,
    rounded corners,
    before skip=10pt,
    after skip=10pt
}

\newtcolorbox{alertaISO}{
    colback=isoAmber,
    colframe=isoOrange,
    boxrule=1pt,
    rounded corners,
    before skip=10pt,
    after skip=10pt
}

%--------------------------------------------------------------------
% TABLAS ANCHO COMPLETO
%--------------------------------------------------------------------
\newenvironment{tablacompleta}
{\begin{adjustwidth}{-0.8cm}{-0.8cm}}
{\end{adjustwidth}}

%--------------------------------------------------------------------
% DROP CAP
%--------------------------------------------------------------------
\usepackage{lettrine}
\lettrinefont{\fontsize{36}{40}\selectfont\bfseries\color{sandraDark}}
\lettrineafterlen=3pt
\lettrinelinenums=false

%--------------------------------------------------------------------
% SEPARADOR
%--------------------------------------------------------------------
\newcommand{\separador}{
    \vspace{1.2ex}
    \begin{center}
        \color{sandraTeal}\rule{0.1\linewidth}{0.8pt}
    \end{center}
    \vspace{1.2ex}
}

%--------------------------------------------------------------------
% PÁGINA DE CAPÍTULO
%--------------------------------------------------------------------
\newcommand{\chapterpage}[2]{
\thispagestyle{empty}
\begin{minipage}[c][\textheight][c]{\linewidth}
    \begin{center}
        \vspace{0.2\textheight}
        {\fontsize{120}{130}\selectfont\bfseries\color{sandraDark!8}#1}
        \vspace{-3cm}
        {\fontsize{30}{36}\selectfont\bfseries\color{sandraDark}#2}
        \vspace{0.5cm}
        {\color{sandraTeal}\rule{0.12\linewidth}{2pt}}
        \vspace{0.3cm}
        {\fontsize{11}{13}\selectfont\color{sandraGray}Capítulo #1}
    \end{center}
\end{minipage}
\clearpage
}

\usepackage[spanish]{babel}

\begin{document}

%======================================================================
% CAPÍTULO 2: MARCO TEÓRICO
%======================================================================
\chapterpage{02}{Marco Teórico}

\section{Estado del Arte}

\lettrine{E}{l} middleware empresarial ha evolucionado significativamente hasta convertirse en plataformas complejas de integración. El mercado ofrece soluciones consolidadas como IBM Integration Bus, MuleSoft, Apache Camel y Spring Integration, cada una con fortalezas específicas para diferentes escenarios de uso.

Sin embargo, la mayoría fueron diseñadas para mercados desarrollados, ignorando las necesidades específicas de regiones con infraestructura limitada y contextos regulatorios distintos. Esta brecha crea una oportunidad clara para soluciones localmente desarrolladas.

\begin{notaeditorial}
    \textbf{Nota del Autor}: El desarrollo de middleware propio representa una oportunidad estratégica para garantizar la soberanía tecnológica en entornos gubernamentales y empresariales críticos.
\end{notaeditorial}

\separador

\section{Fundamentos de Middleware}

El middleware actúa como puente entre aplicaciones y el sistema operativo, facilitando la comunicación e integración entre sistemas heterogéneos. Un Enterprise Service Bus (ESB) representa la evolución hacia una arquitectura basada en servicios.

\begin{estandarISO}
    \textbf{Implementación en Sandra}: Server utiliza Go para proporcionar un ESB de alto rendimiento.
\end{estandarISO}

Las características principales incluyen: transformación de mensajes entre formatos; enrutamiento basado en reglas; composición de servicios; y transacciones distribuidas.

\separador

\section{Arquitectura de Microservicios}

La arquitectura de microservicios construye aplicaciones como servicios pequeños, autónomos y débilmente acoplados. Cada microservicio es responsable de una funcionalidad específica.

\begin{estandarISO}
    \textbf{Ecosistema Sandra}: Server, Consola, SDC, Sentinel
\end{estandarISO}

\begin{tablacompleta}
\begin{table}
    \centering
    \small
    \caption{Componentes del ecosistema}
    \begin{tabular}{lp{7cm}}
        \toprule
        \textbf{Componente} & \textbf{Descripción} \\
        \midrule
        Server & Middleware central de comunicaciones \\
        Consola & Interfaz administrativa web \\
        SDC & Aplicación de escritorio multiplataforma \\
        Sentinel & Procesamiento especializado de eventos \\
        \bottomrule
    \end{tabular}
\end{table}
\end{tablacompleta}

\separador

\section{Modelo Zero Trust}

El modelo Zero Trust elimina la confianza implícita, verificando cada solicitud independientemente de su origen.

\begin{estandarISO}
    \textbf{ISO 27001}: Sistema de Gestión de Seguridad de la Información
\end{estandarISO}

Principios fundamentales: verificación explícita mediante MFA; acceso con mínimo privilegio; verificación continua mediante monitoreo.

\begin{alertaISO}
    \textbf{Consideración}: Zero Trust requiere cambio cultural organizacional.
\end{alertaISO}

\separador

\section{Tecnologías del Ecosistema}

\begin{itemize}
    \item \textbf{Go 1.24}: Rendimiento y concurrencia
    \item \textbf{Rust}: Seguridad de memoria
    \item \textbf{Angular}: Frontend enterprise
    \item \textbf{Tauri}: Aplicaciones de escritorio
    \item \textbf{gRPC}: Comunicación entre servicios
    \item \textbf{WebSockets}: Tiempo real
\end{itemize}

\begin{notaeditorial}
    \textbf{Nota Técnica}: Go + Tauri proporciona equilibrio óptimo.
\end{notaeditorial}

\clearpage

%======================================================================
% CAPÍTULO 3: MARCO LEGAL
%======================================================================
\chapterpage{03}{Marco Legal}

\section{Ley de Infogobierno}

\lettrine{L}{a} Ley de Infogobierno establece normas para el uso de tecnologías en el sector público, promoviendo Soberanía Tecnológica y Software Libre.

\begin{tablacompleta}
\begin{table}
    \centering
    \small
    \caption{Cumplimiento Ley de Infogobierno}
    \begin{tabular}{lp{7cm}}
        \toprule
        \textbf{Requisito} & \textbf{Implementación} \\
        \midrule
        Software Libre & Código abierto \\
        Estándares Abiertos & REST, gRPC, JSON \\
        Interoperabilidad & APIs documentadas \\
        Seguridad & AES-256, JWT, MFA, TLS \\
        \bottomrule
    \end{tabular}
\end{table}
\end{tablacompleta}

\separador

\section{Delitos Informáticos}

Controles implementados:
\begin{itemize}
    \item \textbf{Acceso no autorizado}: MFA
    \item \textbf{Interceptación}: Cifrado
    \item \textbf{Falsificación}: Firmas y logs
    \item \textbf{Destrucción}: Backup
\end{itemize}

\separador

\section{Estándares Internacionales}

\begin{estandarISO}
    \textbf{ISO/IEC 27001}: Gestión de Seguridad
\end{estandarISO}

\begin{estandarISO}
    \textbf{ISO/IEC 27002}: 93 controles
\end{estandarISO}

\begin{estandarISO}
    \textbf{NIST SP 800-53}: Controles de seguridad
\end{estandarISO}

\separador

\section{Licenciamiento Soberano}

\begin{alertaISO}
    \textbf{Nota Estratégica}: Licencias permisivas permiten privatizar innovaciones.
\end{alertaISO}

\begin{tablacompleta}
\begin{table}
    \centering
    \small
    \caption{Análisis de licenciamiento}
    \begin{tabular}{lp{7cm}}
        \toprule
        \textbf{Licencia} & \textbf{Implicación} \\
        \midrule
        Apache 2.0 / MIT & Permite clasificar módulos \\
        GPL v3 / AGPL & Copyleft prohíbe cierre \\
        \bottomrule
    \end{tabular}
\end{table}
\end{tablacompleta}

\begin{notaeditorial}
    \textbf{Nota Final}: El licenciamiento estratégico es crucial.
\end{notaeditorial}

\end{document}
